\begin{figure}[ht]
\centering
\includegraphics[width=\textwidth]{figures/figS15_seed_dependence.pdf}
\caption{\textbf{Supplementary Figure S14. The seed-dependence of prior wiring is explainable by expression gating, not a
methodological failure.}
Each seed's aggregate wiring $z$ (true hematopoietic prior R13 versus the matched-seed empty prior
R14, recomputed here through the same data-quality-gated registry as Supplementary Fig.~S11;
$z_{42}=+2.84$, $z_{123}=+2.01$, $z_{456}=-2.44$ over the $n=317$ prior edges present in the gene
index) is decomposed into a per-TF coupling gain, defined as the
mean over a TF's injected edges of $[\cos(W_{\text{model}}[:,t], W_{\text{model}}[:,g]) -
\cos(W_{\text{vanilla}}[:,t], W_{\text{vanilla}}[:,g])]$ after cross-axis centring and unit
normalisation (identical metric to the wiring test). \textbf{(A)} The sign of the aggregate wiring
$z$ is recovered almost perfectly by the mean per-TF gain across the three seeds (Pearson
$\rho=+0.975$; n=3, a mechanistic decomposition rather than a statistical test). \textbf{(B)}
Per-seed coupling gain of KLF1, the erythroid-core mean, and the all-TF mean: in the two seeds that
wire (42, 123) the gains are positive; in seed~456 the all-TF mean is negative. \textbf{(C)}
Distribution of all 347 prior-edge gains per seed: seeds 42/123 are shifted positive (61\%/55\% of
edges gain), seed~456 is shifted negative (43\% of edges gain). \textbf{Take-home:} the seed~456
flip is idiosyncratic, not a blanket collapse -- in seeds 42/123 the myeloid-core TFs (GATA2, SPI1,
CEBPA) and erythroid TFs (TAL1, GFI1B, LMO2) gain positive coupling, whereas in seed~456 those
\emph{myeloid-core} TFs (SPI1 $-0.149$, CEBPA $-0.083$, NFE2 $-0.209$, GATA2 $-0.106$) flip
negative while the erythroid TFs stay near-neutral. The prior therefore internalises a
seed-stochastic \emph{subset} of its edges -- those whose TFs lie in the trajectory's active program
-- and which lineage program dominates the captured trajectory is set by random initialisation and
cell sampling (expression gating). Seed-dependence is thus a direct, interpretable readout of the
gating mechanism, and the in-silico KLF1 knockout / overexpression in Supplementary Fig.~S12
independently confirms KLF1 as the linchpin of the seed-42 program.}
\label{fig:seed_dependence}
\end{figure}
